\tikzset{
    ddfv/.style={
        scale=1.5,
        every path/.style={thick,line join=round, cap=round},
        ptsq/.style 2 args={
          rectangle,
          fill=##1,
          inner sep=2pt,
          label={[text=##1]##2}
        },
        ptsqbord/.style 2 args={
          rectangle,
          draw=##1,
          inner sep=2pt,
          label={[text=##1]##2}
        },
        ptcl/.style 2 args={
          circle,
          fill=##1,
          inner sep=1.5pt,
          label={[text=##1]##2}
        },
        ptclbord/.style 2 args={
          circle,
          draw=##1,
          inner sep=1.5pt,
          label={[text=##1]##2}
        },
        mailleprimale/.style={
            myblue,
            fill=mylightblue
        },
        mailledualeint/.style={
            myred,
            fill=mylightred
        },
        mailledualebord/.style={
            myred,
            fill=mylightred!50!white
        },
        maillediamantint/.style={
            mygreen,
            fill=mylightgreen
        },
        maillediamantbord/.style={
            mygreen,
            fill=mylightgreen!50!white
        }
    }
}

\def\Rayon{2}
\pgfmathsetmacro{\Rdomaine}{0.9*\Rayon}
\def\xzero{0}
\def\yzero{0}
\def\zzero{0}

\def\eps{0.15}
\def\epsperp{0.05}

\def\exan{1.2}

\newcommand{\MilieuProjete}[2]{%
    % #1 = sommet (A, B, C, ...)
    % #2 = indice

    % Milieu de l'arête primale
    \pgfmathsetmacro{\xM}{(\xzero+\csname x#1\endcsname)/2}
    \pgfmathsetmacro{\yM}{(\yzero+\csname y#1\endcsname)/2}
    \pgfmathsetmacro{\zM}{(f(0,0)+f(\csname x#1\endcsname,\csname y#1\endcsname))/2}

    % Projection radiale sur la sphère de centre (xzero,yzero,zzero)
    \pgfmathsetmacro{\normM}{
        sqrt((\xM-\xzero)^2 + (\yM-\yzero)^2 + (\zM-\zzero)^2)
    }

    \expandafter\pgfmathsetmacro\csname xProjM#1\endcsname
        {\xzero + \Rayon*(\xM-\xzero)/\normM}

    \expandafter\pgfmathsetmacro\csname yProjM#1\endcsname
        {\yzero + \Rayon*(\yM-\yzero)/\normM}

    \expandafter\pgfmathsetmacro\csname zProjM#1\endcsname
        {\zzero + \Rayon*(\zM-\zzero)/\normM}

    \coordinate (fM#1) at (
        {\csname xProjM#1\endcsname},
        {\csname yProjM#1\endcsname},
        {\csname zProjM#1\endcsname}
    );
}

\makeatletter

\newcommand{\PointXYZ}[4]{%
    \pgfmathsetmacro{\px}{#2}%
    \pgfmathsetmacro{\py}{#3}%
    \pgfmathsetmacro{\pz}{#4}%
    \coordinate (#1) at (\px,\py,\pz);
    \expandafter\xdef\csname pt@#1@x\endcsname{\px}%
    \expandafter\xdef\csname pt@#1@y\endcsname{\py}%
    \expandafter\xdef\csname pt@#1@z\endcsname{\pz}%
}

\newcommand{\GetX}[1]{\csname pt@#1@x\endcsname}
\newcommand{\GetY}[1]{\csname pt@#1@y\endcsname}
\newcommand{\GetZ}[1]{\csname pt@#1@z\endcsname}

\newcommand{\CrossPointsFrom}[7]{%
  % #1 = nom du point de départ E (déjà défini)
  % #2 = nom du nouveau point F à créer
  % #3,#4 = A,B   #5,#6 = C,D   #7 = longueur de [E,F]
  \pgfmathsetmacro{\ABx}{\GetX{#4}-\GetX{#3}}%
  \pgfmathsetmacro{\ABy}{\GetY{#4}-\GetY{#3}}%
  \pgfmathsetmacro{\ABz}{\GetZ{#4}-\GetZ{#3}}%
  \pgfmathsetmacro{\CDx}{\GetX{#6}-\GetX{#5}}%
  \pgfmathsetmacro{\CDy}{\GetY{#6}-\GetY{#5}}%
  \pgfmathsetmacro{\CDz}{\GetZ{#6}-\GetZ{#5}}%
  \pgfmathsetmacro{\Nx}{\ABy*\CDz-\ABz*\CDy}%
  \pgfmathsetmacro{\Ny}{\ABz*\CDx-\ABx*\CDz}%
  \pgfmathsetmacro{\Nz}{\ABx*\CDy-\ABy*\CDx}%
  \pgfmathsetmacro{\Nnorm}{sqrt(\Nx*\Nx+\Ny*\Ny+\Nz*\Nz)}%
  \pgfmathsetmacro{\nx}{\Nx/\Nnorm}%
  \pgfmathsetmacro{\ny}{\Ny/\Nnorm}%
  \pgfmathsetmacro{\nz}{\Nz/\Nnorm}%
  \pgfmathsetmacro{\Fx}{\GetX{#1}+#7*\nx}%
  \pgfmathsetmacro{\Fy}{\GetY{#1}+#7*\ny}%
  \pgfmathsetmacro{\Fz}{\GetZ{#1}+#7*\nz}%
  \PointXYZ{#2}{\Fx}{\Fy}{\Fz}%
}

\newcommand{\MidPointXYZ}[3]{%
    \pgfmathsetmacro{\mx}{0.5*(\GetX{#2}+\GetX{#3})}%
    \pgfmathsetmacro{\my}{0.5*(\GetY{#2}+\GetY{#3})}%
    \pgfmathsetmacro{\mz}{0.5*(\GetZ{#2}+\GetZ{#3})}%
    \PointXYZ{#1}{\mx}{\my}{\mz}%
}

\newcommand{\LinePlaneIntersection}[6]{%
    \pgfmathsetmacro{\ux}{\GetX{#5}-\GetX{#4}}%
    \pgfmathsetmacro{\uy}{\GetY{#5}-\GetY{#4}}%
    \pgfmathsetmacro{\uz}{\GetZ{#5}-\GetZ{#4}}%

    \pgfmathsetmacro{\vx}{\GetX{#6}-\GetX{#4}}%
    \pgfmathsetmacro{\vy}{\GetY{#6}-\GetY{#4}}%
    \pgfmathsetmacro{\vz}{\GetZ{#6}-\GetZ{#4}}%

    \pgfmathsetmacro{\nx}{\uy*\vz-\uz*\vy}%
    \pgfmathsetmacro{\ny}{\uz*\vx-\ux*\vz}%
    \pgfmathsetmacro{\nz}{\ux*\vy-\uy*\vx}%

    \pgfmathsetmacro{\dx}{\GetX{#3}-\GetX{#2}}%
    \pgfmathsetmacro{\dy}{\GetY{#3}-\GetY{#2}}%
    \pgfmathsetmacro{\dz}{\GetZ{#3}-\GetZ{#2}}%

    \pgfmathsetmacro{\den}{\nx*\dx+\ny*\dy+\nz*\dz}%
    \pgfmathsetmacro{\num}{\nx*(\GetX{#4}-\GetX{#2})+\ny*(\GetY{#4}-\GetY{#2})+\nz*(\GetZ{#4}-\GetZ{#2})}%

    \pgfmathsetmacro{\tt}{\num/\den}%

    \pgfmathsetmacro{\ix}{\GetX{#2}+\tt*\dx}%
    \pgfmathsetmacro{\iy}{\GetY{#2}+\tt*\dy}%
    \pgfmathsetmacro{\iz}{\GetZ{#2}+\tt*\dz}%

    \PointXYZ{#1}{\ix}{\iy}{\iz}%
}

\makeatother

\pgfmathdeclarefunction{f}{2}{%
  \pgfmathparse{\zzero + sqrt(max(0,\Rayon^2-(#1-\xzero)^2-(#2-\yzero)^2))}%
}

\def\FigureScale{1.8}
\pgfmathsetlengthmacro{\SyncedThickWidth}{0.8pt/\FigureScale}
\tikzset{
    synced thick arrow/.style={
        line width=\SyncedThickWidth,
        -{>[scale={1/\FigureScale}]}
    }
}

\begin{tikzpicture}[
  ddfv,
  scale=\FigureScale,
  node distance=0.6cm,
  every node/.style={inner sep=0pt}
]


\tdplotsetmaincoords{45}{55}
% \tdplotsetmaincoords{15}{-35}

\begin{scope}[
    local bounding box=mainP,
    tdplot_main_coords
]

\coordinate (O) at (0,0,0);
\PointXYZ{fO}{0}{0}{f(0,0)};

\def\Rpenta{1.5}
\def\lengthbasis{4pt}
\def\lengthsq{1pt}

\foreach \i/\j/\angleA/\angleB in {
    1/2/13/17,
    2/3/17/21,
    3/4/21/5,
    4/5/5/9,
    5/1/9/13
}
{
    % Point F_i
    \pgfmathsetmacro{\xA}{\Rpenta*cos((\angleA*pi/10)r)}
    \pgfmathsetmacro{\yA}{\Rpenta*sin((\angleA*pi/10)r)}
    \pgfmathsetmacro{\zA}{f(\xA,\yA)}

    \coordinate (F\i)  at ({\xA},{\yA},0);
    % \coordinate (fF\i) at ({\xA},{\yA},{\zA});
    \PointXYZ{fF\i}{\xA}{\yA}{\zA};

    % \node[ptsq={myblue}{}] at (F\i) {};
    % \node[ptsq={myblue}{}] at (fF\i) {};

    % Point F_j
    \pgfmathsetmacro{\xB}{\Rpenta*cos((\angleB*pi/10)r)}
    \pgfmathsetmacro{\yB}{\Rpenta*sin((\angleB*pi/10)r)}
    \pgfmathsetmacro{\zB}{f(\xB,\yB)}

    \coordinate (F\j)  at ({\xB},{\yB},0);
    \coordinate (fF\j) at ({\xB},{\yB},{\zB});

    % Barycentre euclidien du triangle fO, fF_i, fF_j
    \pgfmathsetmacro{\xG}{(\xzero+\xA+\xB)/3}
    \pgfmathsetmacro{\yG}{(\yzero+\yA+\yB)/3}
    \pgfmathsetmacro{\zG}{(f(0,0)+f(\xA,\yA)+f(\xB,\yB))/3}

    % Projection radiale sur la sphère de centre (xzero,yzero,zzero)
    \pgfmathsetmacro{\normG}{
        sqrt((\xG-\xzero)^2 + (\yG-\yzero)^2 + (\zG-\zzero)^2)
    }

    \pgfmathsetmacro{\xProjG}{\xzero +  \exan*\Rayon*(\xG-\xzero)/\normG}
    \pgfmathsetmacro{\yProjG}{\yzero +  \exan*\Rayon*(\yG-\yzero)/\normG}
    \pgfmathsetmacro{\zProjG}{\zzero +  \exan*\Rayon*(\zG-\zzero)/\normG}

    % \coordinate (fG\i) at ({\xProjG},{\yProjG},{\zProjG});
    \PointXYZ{fG\i}{\xProjG}{\yProjG}{\zProjG};

    %%%%%%%%%%%%%%%%%%%%%%%%%%%%%%%%%%%%%%%%%%%%%%%%%%%%%%%%%%%%%%%%%%%

    \pgfmathsetmacro{\xG}{(\xA+\xB)/3}
    \pgfmathsetmacro{\yG}{(\yA+\yB)/3}
    \pgfmathsetmacro{\zG}{(f(0,0)+f(\xA,\yA)+f(\xB,\yB))/3}

    \PointXYZ{G\i}{\xG}{\yG}{\zG}

    \MidPointXYZ{M\i}{fO}{fF\i}
}

\fill[mylightblue] (fO) -- (fF5) -- (fF1) -- cycle;
\fill[mylightblue] (fO) -- (fF1) -- (fF2) -- cycle;

\draw[mailledualeint] (M1) -- (G1) -- (M2) -- (G2) -- (M3) -- (G3) -- (M4) -- (G4) -- (M5) -- (G5) -- cycle;

\foreach \i in {1,2,3,4,5}
{
    \node[ptcl={myred}{}] at (G\i) {};
}

\foreach \i/\j/\angleA/\angleB in {
    1/2/13/17,
    2/3/17/21,
    3/4/21/5,
    4/5/5/9,
    5/1/9/13
}
{
     % Point F_i
    \pgfmathsetmacro{\xA}{\Rpenta*cos((\angleA*pi/10)r)}
    \pgfmathsetmacro{\yA}{\Rpenta*sin((\angleA*pi/10)r)}
    \pgfmathsetmacro{\zA}{f(\xA,\yA)}

    % Point F_j
    \pgfmathsetmacro{\xB}{\Rpenta*cos((\angleB*pi/10)r)}
    \pgfmathsetmacro{\yB}{\Rpenta*sin((\angleB*pi/10)r)}
    \pgfmathsetmacro{\zB}{f(\xB,\yB)}

    % Barycentre euclidien du triangle fO, fF_i, fF_j
    \pgfmathsetmacro{\xG}{(\xzero+\xA+\xB)/3}
    \pgfmathsetmacro{\yG}{(\yzero+\yA+\yB)/3}
    \pgfmathsetmacro{\zG}{(f(0,0)+f(\xA,\yA)+f(\xB,\yB))/3}

    % Projection radiale sur la sphère de centre (xzero,yzero,zzero)
    \pgfmathsetmacro{\normG}{
        sqrt((\xG-\xzero)^2 + (\yG-\yzero)^2 + (\zG-\zzero)^2)
    }

    \pgfmathsetmacro{\xProjG}{\xzero + \exan*\Rayon*(\xG-\xzero)/\normG}
    \pgfmathsetmacro{\yProjG}{\yzero + \exan*\Rayon*(\yG-\yzero)/\normG}
    \pgfmathsetmacro{\zProjG}{\zzero + \exan*\Rayon*(\zG-\zzero)/\normG}

    %%%%%%%%%%%%%%%%%%%%%%%%%%%%%%%%%%%%%%%%%%%%%%%%%%%%%%%%%%%%%%%%%%%

    \pgfmathsetmacro{\xG}{(\xA+\xB)/3}
    \pgfmathsetmacro{\yG}{(\yA+\yB)/3}
    \pgfmathsetmacro{\zG}{(f(0,0)+f(\xA,\yA)+f(\xB,\yB))/3}

    \PointXYZ{G\i}{\xG}{\yG}{\zG}

    % \node[ptcl={myred}{}] at (G\i) {};
    % \draw[myred,dotted] (fG\i) -- (G\i);

    % Normale radiale unitaire en G_i
    \pgfmathsetmacro{\nx}{(\xProjG-\xzero)/\Rayon}
    \pgfmathsetmacro{\ny}{(\yProjG-\yzero)/\Rayon}
    \pgfmathsetmacro{\nz}{(\zProjG-\zzero)/\Rayon}

    % Première direction tangente : e_theta = (-ny,nx,0)
    \pgfmathsetmacro{\nt}{sqrt(\nx*\nx+\ny*\ny)}
    \pgfmathsetmacro{\tx}{-\ny/\nt}
    \pgfmathsetmacro{\ty}{ \nx/\nt}
    \pgfmathsetmacro{\tz}{0}

    % Deuxième direction tangente : s = n x t
    \pgfmathsetmacro{\sx}{-\nz*\ty}
    \pgfmathsetmacro{\sy}{ \nz*\tx}
    \pgfmathsetmacro{\sz}{ \nx*\ty-\ny*\tx}


    \onslide<2->{
    % Équerre dans le plan (n,t)
    \draw[semithick]
    (G\i)
    -- ++({\epsperp*\tx},{\epsperp*\ty},{\epsperp*\tz})
    -- ++({\epsperp*\nx},{\epsperp*\ny},{\epsperp*\nz})
    -- ++({-\epsperp*\tx},{-\epsperp*\ty},{-\epsperp*\tz})
    -- cycle;

    % Équerre dans le plan (n,s)
    \draw[semithick]
    (G\i)
    -- ++({\epsperp*\sx},{\epsperp*\sy},{\epsperp*\sz})
    -- ++({\epsperp*\nx},{\epsperp*\ny},{\epsperp*\nz})
    -- ++({-\epsperp*\sx},{-\epsperp*\sy},{-\epsperp*\sz})
    -- cycle;
    }

    %%%%%%%%%%%%%%%%%%%%%%%%%%%%%%%%%%%%%%%%%%%%%%%%%%%%%%%%%%%%%%%%%%%

    % \draw[myorange] (fO) -- (G\i) -- (fF\i);
    % \draw[myorange] (G\i) -- (fF\j);
}

% \ruledsurface{fO}{fG2}{fF3}{fG3}
    
% \node[myblue, above=6pt] at (fO) {\contour{mylightred!50!white}{$x_{\mailledualeplate{K}}$}};
% \node[above=5pt] at (M2) {\contour{mylightred!50!white}{$x_\sigma$}};
% \node[myblue] at ($(fO)!0.8!(fF2)$) {\contour{white}{$\sigma_{\primalplat{K},\primalplat{L}}$}};
% \node[myred, below=5pt] at (G1) {$x_{\primalplat{L}}$};
% \node[myred, above=5pt, visible on=<2->] at (fG1) {\contour{white}{$x_{\primal{L}}$}};
% \node[myred, below=5pt] at (G2) {$x_{\primalplat{K}}$};
% \node[myred, above=5pt, visible on=<2->] at (fG2) {$x_{\primal{K}}$};
% \node[myblue, left=5pt] at ($(fF1)!0.5!(fF2)$) {$\primalplat{L}$};
% \node[myblue, below=5pt] at ($(fF2)!0.5!(fF3)$) {$\primalplat{K}$};
% \node[myred, below=5pt] at ($(G2)!0.5!(M2)$) {$\mailledualeplate{K}$};

\CrossPointsFrom{fF1}{fF1v}{fF1}{fO}{fG5}{fG1}{1.5}
\CrossPointsFrom{fO}{fOv}{fF1}{fO}{fG5}{fG1}{1.5}
\CrossPointsFrom{fG5}{fG5v}{fF1}{fO}{fG5}{fG1}{1.5}
\CrossPointsFrom{fG1}{fG1v}{fF1}{fO}{fG5}{fG1}{1.5}

\CrossPointsFrom{fF2}{fF2v}{fF2}{fO}{fG1}{fG2}{1.5}
\CrossPointsFrom{fO}{fOv}{fF2}{fO}{fG1}{fG2}{1.5}
\CrossPointsFrom{fG1}{fG1v}{fF2}{fO}{fG1}{fG2}{1.5}
\CrossPointsFrom{fG2}{fG2v}{fF2}{fO}{fG1}{fG2}{1.5}

\MidPointXYZ{fF1G1}{fF1}{fG1}
\MidPointXYZ{fOG1}{fO}{fG1}
\MidPointXYZ{fOG5}{fO}{fG5}
\MidPointXYZ{fF1G5}{fF1}{fG5}

% Plan du premier diamant : il passe par le barycentre des quatre
% sommets et contient les directions fO--fF1 et fG5--fG1.
\PointXYZ{DiamondFirstO}
    {(\GetX{fO}+\GetX{fF1}+\GetX{fG5}+\GetX{fG1})/4}
    {(\GetY{fO}+\GetY{fF1}+\GetY{fG5}+\GetY{fG1})/4}
    {(\GetZ{fO}+\GetZ{fF1}+\GetZ{fG5}+\GetZ{fG1})/4}
\PointXYZ{DiamondFirstU}
    {\GetX{DiamondFirstO}+\GetX{fF1}-\GetX{fO}}
    {\GetY{DiamondFirstO}+\GetY{fF1}-\GetY{fO}}
    {\GetZ{DiamondFirstO}+\GetZ{fF1}-\GetZ{fO}}
\PointXYZ{DiamondFirstV}
    {\GetX{DiamondFirstO}+\GetX{fG1}-\GetX{fG5}}
    {\GetY{DiamondFirstO}+\GetY{fG1}-\GetY{fG5}}
    {\GetZ{DiamondFirstO}+\GetZ{fG1}-\GetZ{fG5}}

\MidPointXYZ{fF2G2}{fF2}{fG2}
\MidPointXYZ{fOG2}{fO}{fG2}
\MidPointXYZ{fOG1}{fO}{fG1}
\MidPointXYZ{fF2G1}{fF2}{fG1}

% \node[ptsq={black}{}] at (fF1G1) {};
% \node[ptsq={black}{}] at (fOG1) {};
% \node[ptsq={black}{}] at (fOG5) {};
% \node[ptsq={black}{}] at (fF1G5) {};

\LinePlaneIntersection{I}{fF1}{fF1v}{DiamondFirstO}{DiamondFirstU}{DiamondFirstV}
\LinePlaneIntersection{K}{fG5}{fG5v}{DiamondFirstO}{DiamondFirstU}{DiamondFirstV}

% I et K fixent la position du plan. J et L sont ensuite obtenus par
% translation afin de conserver exactement les deux directions voulues.
\PointXYZ{J}
    {\GetX{I}+\GetX{fO}-\GetX{fF1}}
    {\GetY{I}+\GetY{fO}-\GetY{fF1}}
    {\GetZ{I}+\GetZ{fO}-\GetZ{fF1}}
\PointXYZ{L}
    {\GetX{K}+\GetX{fG1}-\GetX{fG5}}
    {\GetY{K}+\GetY{fG1}-\GetY{fG5}}
    {\GetZ{K}+\GetZ{fG1}-\GetZ{fG5}}

\LinePlaneIntersection{Ib}{fF2}{fF2v}{fF2G1}{fOG1}{fOG2}
\LinePlaneIntersection{Jb}{fO}{fOv}{fF2G1}{fOG1}{fOG2}
\LinePlaneIntersection{Kb}{fG1}{fG1v}{fF2G2}{fOG1}{fOG2}
\LinePlaneIntersection{Lb}{fG2}{fG2v}{fF2G1}{fOG1}{fOG2}

\draw[myblue] (fF1) -- (fF2) -- (fF3) -- (fF4) -- (fF5) -- cycle;

\LinePlaneIntersection{InterG5}{G5}{fG5}{J}{K}{I}
\LinePlaneIntersection{InterG1}{G1}{fG1}{J}{K}{I}

\foreach \i in {1,2,3,4,5}
{
    % \draw (G\i) -- (fO);
    \draw[myblue] (fO) -- (fF\i);
    \draw[dashed, thin, visible on=<2->] (G\i) -- (fG\i);
    \node[ptsq={black}{}] at (M\i) {};
}

\onslide<4->{

\draw[black] (fO) -- (fG2) -- (fF3);
\draw[black] (fO) -- (fG3) -- (fF3);

}

\node[ptsq={myblue}{}] at (fO) {};

\foreach \i in {1,2,3,4,5}
{
    \node[ptsq={myblue}{}] at (fF\i) {};
}

% Couleurs de sphere_diamond_barycentric_gmsh.py (noms CSS/PyVista).
\definecolor{diamondPythonGreen}{HTML}{008000}
\definecolor{diamondPythonLimeGreen}{HTML}{32CD32}
\filldraw[fill=diamondPythonGreen, draw=diamondPythonLimeGreen, opacity=0.8, visible on=<3->] (J) -- (K) -- (I) -- (L) -- cycle;

\node[diamondPythonGreen, align=center, font=\footnotesize, visible on=<3->] at ($(fF1G5)+(180:8pt)$) {diamant\\plat};

% \filldraw[mygreen, draw=mydarkgreen, opacity=0.8, visible on=< 3->] (J) -- (K) -- (I) -- (L) -- cycle;
% \draw[maillediamantint, opacity=0.8, thin, visible on=<4->] (Jb) -- (Kb) -- (Ib) -- (Lb) -- cycle;

\onslide<3->{
% Repère affine du plan du diamant :
%   origine Jb,
%   axe x porté par Jb--Ib,
%   axe y parallèle à Kb--Lb.
\PointXYZ{DiamondPlaneY}
    {\GetX{Jb}+\GetX{Lb}-\GetX{Kb}}
    {\GetY{Jb}+\GetY{Lb}-\GetY{Kb}}
    {\GetZ{Jb}+\GetZ{Lb}-\GetZ{Kb}}

% Coordonnées de Kb dans la base non nécessairement orthogonale
% u = Ib-Jb et v = Lb-Kb.
\pgfmathsetmacro{\dux}{\GetX{Ib}-\GetX{Jb}}
\pgfmathsetmacro{\duy}{\GetY{Ib}-\GetY{Jb}}
\pgfmathsetmacro{\duz}{\GetZ{Ib}-\GetZ{Jb}}
\pgfmathsetmacro{\dvx}{\GetX{Lb}-\GetX{Kb}}
\pgfmathsetmacro{\dvy}{\GetY{Lb}-\GetY{Kb}}
\pgfmathsetmacro{\dvz}{\GetZ{Lb}-\GetZ{Kb}}
\pgfmathsetmacro{\dwx}{\GetX{Kb}-\GetX{Jb}}
\pgfmathsetmacro{\dwy}{\GetY{Kb}-\GetY{Jb}}
\pgfmathsetmacro{\dwz}{\GetZ{Kb}-\GetZ{Jb}}

\pgfmathsetmacro{\duu}{\dux*\dux+\duy*\duy+\duz*\duz}
\pgfmathsetmacro{\duv}{\dux*\dvx+\duy*\dvy+\duz*\dvz}
\pgfmathsetmacro{\dvv}{\dvx*\dvx+\dvy*\dvy+\dvz*\dvz}
\pgfmathsetmacro{\dwu}{\dwx*\dux+\dwy*\duy+\dwz*\duz}
\pgfmathsetmacro{\dwv}{\dwx*\dvx+\dwy*\dvy+\dwz*\dvz}
\pgfmathsetmacro{\ddet}{\duu*\dvv-\duv*\duv}
\pgfmathsetmacro{\Klocalx}{(\dwu*\dvv-\dwv*\duv)/\ddet}
\pgfmathsetmacro{\Klocaly}{(\dwv*\duu-\dwu*\duv)/\ddet}

\begin{scope}[
    plane origin={(Jb)},
    plane x={(Ib)},
    plane y={(DiamondPlaneY)},
    canvas is plane,
]

\coordinate (xle2D) at (0,0);                % Jb
\coordinate (xke2D) at (1,0);                % Ib
\coordinate (xk2D)  at (\Klocalx,\Klocaly); % Kb
\coordinate (xl2D)  at (\Klocalx,{\Klocaly+1}); % Lb

% \fill[mylightgreen,opacity=0.8] (xle2D) -- (xk2D) -- (xke2D) -- (xl2D) -- cycle;

% \draw[maillediamantint, opacity=0.8, thin, visible on=<4->] (Jb) -- (Kb) -- (Ib) -- (Lb) -- cycle;

\path[name path=sigma]  (xle2D) -- (xke2D);
\path[name path=sigmae] (xk2D) -- (xl2D);

\path[name intersections={of=sigma and sigmae, by=Inter}];

% \draw[mygreen!80!white, thin] (InterG1) -- (fG1);
\draw[dashed, thin] (InterG1) -- (fG1);

\draw[white, thin] (InterG5) -- (fG5);
\draw[dashed, thin] (InterG5) -- (fG5);

\iffalse
\def\rayonangle{0.08}

\path let
    \p0 = (Inter),
    \p1 = (xke2D),
    \p2 = (xl2D)
in
    \pgfextra{
        \pgfmathsetmacro{\angA}{atan2(-\y1+\y0,\x1-\x0)}
        \pgfmathsetmacro{\angB}{atan2(\y2-\y0,\x2-\x0)}
        \global\let\angA\angA
        \global\let\angB\angB
    };

\draw[black]
    ($(Inter)+(\angA:\rayonangle)$)
    arc[
        start angle=\angA,
        end angle=\angB,
        radius=\rayonangle
    ];

\node at ($(Inter)+({0.5*(\angA+\angB)}:{2*\rayonangle})$) {$\angle{\diamantplat{D}}$};

\draw ($(xke2D)!0.2!(xle2D)$) rectangle ++(\lengthsq,-\lengthsq);
\coordinate (P) at ($(xk2D)!0.2!(xl2D)$);
\draw
    (P)
    -- ($(P)!\lengthsq!(xl2D)$)
    -- ($(P)!\lengthsq!(xl2D)!\lengthsq!-90:(xl2D)$)
    -- ($(P)!\lengthsq!-90:(xl2D)$)
    -- cycle;


\coordinate (Pbase) at ($(xke2D)!0.2!(xle2D)$);
\coordinate (Pbasedual) at ($(xk2D)!0.2!(xl2D)$);

\node[myblue, above left] at ($(Pbase)+(0,{-\lengthbasis})$) {$\vect{\tau}_{\mailledualeplate{K},\mailledualeplate{L}}$};
\node[myblue, right] at ($(Pbase)+({\lengthbasis},0)$) {$\vect{\nu}_{\sigma \primalplat{K}}$};
    
\node[myred, above] at ($(Pbasedual)!\lengthbasis!(xl2D)$) {$\vect{\tau}_{\primalplat{K},\primalplat{L}}$};
\node[myred, below] at ($(Pbasedual)!\lengthbasis!-90:(xl2D)$) {$\vect{\nu}_{\dual{\sigma}\mailledualeplate{K}}$};
\fi

\end{scope}

% \draw[myblue, dashed] (Ib) -- (Jb);
% \draw[myred, dashed] (Kb) -- (Lb);

\draw[myblue, dashed] (I) -- (J);
\draw[myred, dashed] (K) -- (L);
}

\iffalse
\onslide<4->{

\draw[mygreen!80!white, thin] (InterG1) -- (fG1);
\draw[dashed, thin] (InterG1) -- (fG1);

\draw[white, thin] (InterG5) -- (fG5);
\draw[dashed, thin] (InterG5) -- (fG5);

\draw[gray] (fOG2) -- (fG2);
\draw[gray] (fF2G2) -- (fG2);
\draw[gray] (fOG1) -- (fG1);
\draw[gray] (fF2G1) -- (fG1);
\draw[gray] (fF1G1) -- (fG1);

% \node[mydarkgreen] at ($(Inter)+(120:8pt)$) {$\diamantplat{D}$};
}
\fi

\onslide<2->{
    \draw[myred] (fG1) -- (fG5);
}


\onslide<5->{
%%%%%%%%%%%%%%%%%%%%%%%%%%%%%%%%%%%%%%%%%%%%%%%%%%%%%%%%%%%%
%                    Diamant courbe
%%%%%%%%%%%%%%%%%%%%%%%%%%%%%%%%%%%%%%%%%%%%%%%%%%%%%%%%%%%%
\foreach \i/\j/\k/\angleA/\angleB/\angleC in {
    3/4/2/21/5/17 % 1/2/5/13/17/9
}
{
    % Point F_i
    \pgfmathsetmacro{\xA}{\Rpenta*cos((\angleA*pi/10)r)}
    \pgfmathsetmacro{\yA}{\Rpenta*sin((\angleA*pi/10)r)}
    \pgfmathsetmacro{\zA}{f(\xA,\yA)}

    % Point F_j
    \pgfmathsetmacro{\xB}{\Rpenta*cos((\angleB*pi/10)r)}
    \pgfmathsetmacro{\yB}{\Rpenta*sin((\angleB*pi/10)r)}
    \pgfmathsetmacro{\zB}{f(\xB,\yB)}

    % Point F_k
    \pgfmathsetmacro{\xC}{\Rpenta*cos((\angleC*pi/10)r)}
    \pgfmathsetmacro{\yC}{\Rpenta*sin((\angleC*pi/10)r)}
    \pgfmathsetmacro{\zC}{f(\xC,\yC)}

    % Barycentre euclidien du triangle fO, fF_i, fF_j
    \pgfmathsetmacro{\xG}{(\xzero+\xA+\xB)/3}
    \pgfmathsetmacro{\yG}{(\yzero+\yA+\yB)/3}
    \pgfmathsetmacro{\zG}{(f(0,0)+f(\xA,\yA)+f(\xB,\yB))/3}

    % Projection radiale sur la sphère de centre (xzero,yzero,zzero)
    \pgfmathsetmacro{\normG}{
        sqrt((\xG-\xzero)^2 + (\yG-\yzero)^2 + (\zG-\zzero)^2)
    }

    \pgfmathsetmacro{\xProjG}{\xzero + \exan*\Rayon*(\xG-\xzero)/\normG}
    \pgfmathsetmacro{\yProjG}{\yzero + \exan*\Rayon*(\yG-\yzero)/\normG}
    \pgfmathsetmacro{\zProjG}{\zzero + \exan*\Rayon*(\zG-\zzero)/\normG}

    
    % Barycentre euclidien du triangle fO, fF_i, fF_k
    \pgfmathsetmacro{\xGb}{(\xzero+\xA+\xC)/3}
    \pgfmathsetmacro{\yGb}{(\yzero+\yA+\yC)/3}
    \pgfmathsetmacro{\zGb}{(f(0,0)+f(\xA,\yA)+f(\xC,\yC))/3}

    % Projection radiale sur la sphère de centre (xzero,yzero,zzero)
    \pgfmathsetmacro{\normGb}{
        sqrt((\xGb-\xzero)^2 + (\yGb-\yzero)^2 + (\zGb-\zzero)^2)
    }

    \pgfmathsetmacro{\xProjGb}{\xzero + \exan*\Rayon*(\xGb-\xzero)/\normGb}
    \pgfmathsetmacro{\yProjGb}{\yzero + \exan*\Rayon*(\yGb-\yzero)/\normGb}
    \pgfmathsetmacro{\zProjGb}{\zzero + \exan*\Rayon*(\zGb-\zzero)/\normGb}

    % Écarts entre les points fG imposés et la sphère de référence.
    % La correction est maximale au sommet rouge et s'annule aux
    % extrémités bleues du diamant.
    \pgfmathsetmacro{\DeltaZG}{\zProjG-f(\xProjG,\yProjG)}
    \pgfmathsetmacro{\DeltaZGb}{\zProjGb-f(\xProjGb,\yProjGb)}

    \definecolor{light_viridian}{HTML}{6EFF70}
    \fill[light_viridian, opacity=0.8]
    plot[variable=\t, domain=1:0, samples=20]
    (
        {\t*\xProjG},
        {\t*\yProjG},
        {f(\t*\xProjG, \t*\yProjG)+\t*\DeltaZG}
    )
    --
    plot[variable=\t, domain=1:0, samples=20]
    (
        {(1-\t)*\xProjG + \t*\xA},
        {(1-\t)*\yProjG + \t*\yA},
        {f((1-\t)*\xProjG + \t*\xA, (1-\t)*\yProjG + \t*\yA)
            +(1-\t)*\DeltaZG}
    ) 
    --
    plot[variable=\t, domain=0:1, samples=20]
    (
        {\t*\xProjGb},
        {\t*\yProjGb},
        {f(\t*\xProjGb, \t*\yProjGb)+\t*\DeltaZGb}
    )
    --
    plot[variable=\t, domain=0:1, samples=20]
    (
        {(1-\t)*\xProjGb + \t*\xA},
        {(1-\t)*\yProjGb + \t*\yA},
        {f((1-\t)*\xProjGb + \t*\xA, (1-\t)*\yProjGb + \t*\yA)
            +(1-\t)*\DeltaZGb}
    )
    --cycle;

    % Les quatre bords sont tracés comme des courbes ouvertes : les
    % fermetures droites nécessaires au remplissage restent invisibles.
    \draw[diamondPythonGreen, thick]
    plot[variable=\t, domain=0:1, samples=20]
    (
        {\t*\xProjG},
        {\t*\yProjG},
        {f(\t*\xProjG, \t*\yProjG)+\t*\DeltaZG}
    );

    \draw[diamondPythonGreen, thick]
    plot[variable=\t, domain=0:1, samples=20]
    (
        {(1-\t)*\xProjG + \t*\xA},
        {(1-\t)*\yProjG + \t*\yA},
        {f((1-\t)*\xProjG + \t*\xA, (1-\t)*\yProjG + \t*\yA)
            +(1-\t)*\DeltaZG}
    );

    \draw[diamondPythonGreen, thick]
    plot[variable=\t, domain=1:0, samples=20]
    (
        {\t*\xProjGb},
        {\t*\yProjGb},
        {f(\t*\xProjGb, \t*\yProjGb)+\t*\DeltaZGb}
    );

    \draw[diamondPythonGreen, thick]
    plot[variable=\t, domain=1:0, samples=20]
    (
        {(1-\t)*\xProjGb + \t*\xA},
        {(1-\t)*\yProjGb + \t*\yA},
        {f((1-\t)*\xProjGb + \t*\xA, (1-\t)*\yProjGb + \t*\yA)
            +(1-\t)*\DeltaZGb}
    );
}
%%%%%%%%%%%%%%%%%%%%%%%%%%%%%%%%%%%%%%%%%%%%%%%%%%%%%%%%%%%%
%                    fin Diamant courbe
%%%%%%%%%%%%%%%%%%%%%%%%%%%%%%%%%%%%%%%%%%%%%%%%%%%%%%%%%%%%
}

\node[light_viridian!70!black, align=center, font=\footnotesize, visible on=<5->] at ($(G2)+(-25:8pt)$) {diamant\\courbe};

% \begin{scope}[canvas is xz plane at y=0]
% \path  (0,1) node[anchor=south]{$z$}; 
% \pgflowlevelsynccm
% \draw[thick,-latex] (0,0) -- (0,1); 
% \end{scope}

% Repère affine du premier diamant : axe x porté par J--I et
% axe y parallèle à K--L.
\PointXYZ{DiamondPlaneYFirst}
    {\GetX{Jb}+\GetX{Lb}-\GetX{Kb}}
    {\GetY{Jb}+\GetY{Lb}-\GetY{Kb}}
    {\GetZ{Jb}+\GetZ{Lb}-\GetZ{Kb}}

\pgfmathsetmacro{\fdux}{\GetX{Ib}-\GetX{Jb}}
\pgfmathsetmacro{\fduy}{\GetY{Ib}-\GetY{Jb}}
\pgfmathsetmacro{\fduz}{\GetZ{Ib}-\GetZ{Jb}}
\pgfmathsetmacro{\fdvx}{\GetX{Lb}-\GetX{Kb}}
\pgfmathsetmacro{\fdvy}{\GetY{Lb}-\GetY{Kb}}
\pgfmathsetmacro{\fdvz}{\GetZ{Lb}-\GetZ{Kb}}
\pgfmathsetmacro{\fdwx}{\GetX{Kb}-\GetX{Jb}}
\pgfmathsetmacro{\fdwy}{\GetY{Kb}-\GetY{Jb}}
\pgfmathsetmacro{\fdwz}{\GetZ{Kb}-\GetZ{Jb}}

\pgfmathsetmacro{\fduu}{\fdux*\fdux+\fduy*\fduy+\fduz*\fduz}
\pgfmathsetmacro{\fduv}{\fdux*\fdvx+\fduy*\fdvy+\fduz*\fdvz}
\pgfmathsetmacro{\fdvv}{\fdvx*\fdvx+\fdvy*\fdvy+\fdvz*\fdvz}
\pgfmathsetmacro{\fdwu}{\fdwx*\fdux+\fdwy*\fduy+\fdwz*\fduz}
\pgfmathsetmacro{\fdwv}{\fdwx*\fdvx+\fdwy*\fdvy+\fdwz*\fdvz}
\pgfmathsetmacro{\fddet}{\fduu*\fdvv-\fduv*\fduv}
\pgfmathsetmacro{\FirstKlocalx}{(\fdwu*\fdvv-\fdwv*\fduv)/\fddet}
\pgfmathsetmacro{\FirstKlocaly}{(\fdwv*\fduu-\fdwu*\fduv)/\fddet}

\iffalse
\begin{scope}[
    plane origin={(Jb)},
    plane x={(Ib)},
    plane y={(DiamondPlaneYFirst)},
    canvas is plane,
]
    \pgflowlevelsynccm

    \coordinate (xleZero) at (0,0); % Jb
    \coordinate (xkeZero) at (1,0); % Ib
    \coordinate (xkZero) at (\FirstKlocalx,\FirstKlocaly); % Kb
    \coordinate (xlZero) at (\FirstKlocalx,{\FirstKlocaly+1}); % Lb

    \coordinate (Pbase) at ($(xkeZero)!0.2!(xleZero)$);
    \coordinate (Pbasedual) at ($(xkZero)!0.2!(xlZero)$);

    \draw[myblue, synced thick arrow] (Pbase) -- ++(0,-\lengthbasis); % node[above=5pt] {$\vect{\tau}_{\mailleduale{K},\mailleduale{L}}$};
    \draw[myblue, synced thick arrow] (Pbase) -- ++(\lengthbasis,0); % node[above] {$\vect{\nu}_{\sigma \primal{K}}$};

    \draw[myred, synced thick arrow] (Pbasedual) -- ($(Pbasedual)!\lengthbasis!(xlZero)$);
    \draw[myred, synced thick arrow] (Pbasedual) -- ($(Pbasedual)!\lengthbasis!-90:(xlZero)$); % node[below] {$\vect{\nu}_{\dual{\sigma}\mailleduale{K}}$};

\end{scope}
\fi

\foreach \i in {1,2,3,4,5}
{
    \node[ptcl={myred}{}, visible on=<2->] at (fG\i) {};
}

\end{scope}

\end{tikzpicture}
